\chapter{Kết luận và khuyến nghị}

\section{Kết luận chung}
Báo cáo đã hoàn thành quy trình rút trích dữ liệu và tiền xử lý cho bài toán
dự đoán tiểu đường trên bộ dữ liệu Pima. Kết quả cho thấy việc kết hợp
\textit{zero-as-missing} với pipeline imputation + scaling là cần thiết để nâng cao
tính nhất quán của dữ liệu và hiệu năng mô hình baseline.

\section{Tóm tắt kết quả đạt được}
\begin{itemize}[leftmargin=1.5em]
    \item Hoàn thiện notebook phân tích theo hướng báo cáo học thuật.
    \item So sánh 12 cấu hình tiền xử lý và mô hình trên cùng điều kiện chia dữ liệu.
    \item Đề xuất cấu hình KNN + MinMax + Random Forest cho kết quả test ROC-AUC cao nhất.
    \item Bổ sung phân tích confusion matrix, ROC và threshold trade-off.
\end{itemize}

\section{Trả lời trực tiếp câu hỏi nghiên cứu (Q1/Q2/Q3)}
\begin{itemize}[leftmargin=1.5em]
    \item \textbf{Q1 -- Zero-as-missing có hữu ích không?}
    Có. Việc chuyển \texttt{0 -> NaN} cho các biến sinh lý giúp mô hình đạt mức AUC ổn định
    (test ROC-AUC quanh 0.81) thay vì để nhiễu từ giá trị 0 phi sinh lý.
    \item \textbf{Q2 -- Tổ hợp preprocessing/model nào tốt nhất?}
    Cấu hình tốt nhất là \textbf{KNN + MinMax + Random Forest} với
    \texttt{val\_roc\_auc=0.8700}, \texttt{test\_roc\_auc=0.8170}, \texttt{test\_accuracy=0.7403}.
    \item \textbf{Q3 -- Trade-off precision/recall thay đổi ra sao theo threshold?}
    Khi giảm ngưỡng từ 0.50 xuống 0.30, recall tăng mạnh từ 0.5741 lên 0.8148, trong khi precision
    giảm từ 0.6458 xuống 0.5789; F1 cải thiện từ 0.6078 lên 0.6769.
\end{itemize}

\section{Đóng góp của đề tài}
\begin{itemize}[leftmargin=1.5em]
    \item Chuẩn hóa quy trình thực hành từ EDA đến đánh giá mô hình baseline.
    \item Cung cấp khung đánh giá dễ tái lập cho các bài toán phân lớp y sinh tương tự.
    \item Tăng tính báo cáo hóa của notebook để phục vụ trình bày học phần.
\end{itemize}

\section{Hạn chế nghiên cứu}
\begin{enumerate}[leftmargin=1.8em]
    \item Quy mô dữ liệu còn nhỏ (768 mẫu), chưa đại diện cho nhiều quần thể bệnh nhân.
    \item Giả định \textit{zero-as-missing} có thể gây sai lệch nếu tồn tại trường hợp đo thực sự bằng 0.
    \item Đánh giá hiện tại dựa trên một lần chia train/validation/test, chưa có CV lặp lại.
    \item Chưa tiến hành tuning tham số hệ thống và calibration xác suất.
\end{enumerate}

\section{Hướng phát triển và kiến nghị}
\begin{enumerate}[leftmargin=1.8em]
    \item Áp dụng Stratified K-Fold CV để kiểm chứng độ ổn định kết quả.
    \item Mở rộng tuning tham số cho Logistic Regression và Random Forest.
    \item Thử nghiệm thêm các mô hình boosting (XGBoost/LightGBM) và kỹ thuật cân bằng lớp.
    \item Mở rộng bộ dữ liệu và kiểm tra external validity trên nguồn độc lập.
    \item Đóng gói pipeline suy luận và tạo báo cáo tự động để tăng khả năng triển khai.
\end{enumerate}
